Clinical CT datasets are rarely annotated for a single target: a classification cohort carries many findings and a segmentation dataset many structures. Foundation models for 3-D CT are meant to be adapted to such datasets, yet their published evaluations mostly keep the encoder frozen and train a light head on its features, since this is cheap and measures the pretrained representation rather than the fine-tuning procedure. However, adapting a model to a multi-label cohort or a multi-structure dataset requires two decisions that frozen evaluations do not inform: which parameters to train, and whether to train one model per label or one model for the whole dataset. In this work, we go beyond frozen encoders and compare \numClsMethodsN{} classification and \numSegMethodsN{} segmentation adaptation methods, from frozen probes to low-rank adapters and full fine-tuning, each trained per label and jointly, in an effort to find out what to train and when one model can serve all labels of a dataset. We consider four public 3-D CT encoders and use \numLabelsN{} labels from five public classification cohorts and \numStructuresN{} structures from five public segmentation datasets (\numRunsMainN{} runs). In classification, adapting the attention weights of the encoder improved on frozen probes on three of the four encoders, and one joint model per cohort was on average level with per-label models. In segmentation, a multi-depth decoder on the frozen encoder gave most of the gain, while one model per structure remained more accurate than a joint model at our training budgets. We condense the findings into a practical guide and release code, splits and per-case predictions.
